\cotefile{72}
\scanpage{82}

% L1–2, top
\[
\text{donc}\quad
\underbrace{\operatorname{lin}(h_{i+1},\ldots,h_n)}_{Z_i\ \dim n-i\unsure{-1|+1}}
 \subset Z\subset
\underbrace{H_0\cap\cdots\cap H_i}_{Z'_i\ \dim n-i}.
\]

% L3–5
\deleted{Celà implique donc $Z$} et que cette condition implique bien que $Z$ est stable \unsure{par|pour} les $\unsure{\tau_j|\mathrm T_j}$. Elle \deleted{\unsure{ex}}implique aussi

% L6
\[
\deleted{Z.}\qquad Z_i\subset Z'_i
\]

% L7–9
ce qui équivaut aussi à

% L8–9, centre
\[
\begin{aligned}
 h_{i+1}&\notin H_i\\
 \text{i.e.}\qquad \beta_i&=0
\end{aligned}
\]

% L6–10, right
\margin{Dégager, pour $i$ fixé conditions équiv.:\\
a) $Z_i\subset Z'_i$\\
b) $Z_i$ inv. par $G$\\
c) $Z'_i$ inv. par $G$\\
d) $\beta_i=0$}

% L10–11
Et par une raison de dim, on doit avoir

% L11
\[
Z=Z_i\ \text{ou}\ Z=Z'_i.
\]

% L12
L'inverse étant clair.

% L13–16
\uline{Remarque} Conditions équivalentes, lorsque $\dim C\geqslant 2$ i.e $n\geqslant 2$ : a) Il n'y a pas d'hyperplan de $C$ invariant par $G$ b) On \deleted{\unsure{a}} a \deleted{\unsure{pas}} $\beta_0\ne0$, et \deleted{\unsure{est}} \inserted{on n'a pas}

% L16–19, underbrace and interline
\[
\underbrace{\left(\beta_1\deleted{\ne}=0\
 \deleted{\text{et}}\ 
 \overset{\inserted{\text{(i.e. }H_0\cap H_1\ \text{inv}\deleted{\unsure{.}}\ \deleted{\unsure{\text{ou}}}\ \text{et)}}}
 {\underbrace{H_0\cap H_1}_{\dim n-1}\subset C}\right)}
 _{\substack{\Updownarrow\ \text{si }\beta_0\ \text{ou }\mu_0\ne0\\
 \beta_1=0\ \text{et }\mu_1\mu_0-\beta_0=0}}.
\]

% archive sheet number: 81
\scanpage{83}

% L1
espaces forment un drapeau.

% L2–6
e) Pour que $G$ opérant sur $C$ soit irréductible, il f et s. que $\rho$ et les $\beta_i$ ($0\leq i\leq n-1$) soient $\ne0$, i.e. que \deleted{\unsure{$f$}} la forme \deleted{quadratique can. $f_V$} \inserted{bil. can. $\unsure{\varphi_V|\psi_V|f_V}$} sur $V$ soit \deleted{non dég.} non dégénérée (i.e. que $f_V$ soit lisse, et que \deleted{\unsure{pour}} \inserted{si} $n$ impair, \deleted{$2$ soit inv. en $S$} \inserted{$k$ soit de car. $\ne2$})

% L7
\deleted{\unsure{f)|c)} Pour que}

% L8–11
\uline{Dém} Dire que $Z$ est inv. par $\unsure{\tau_i|\mathrm T_i}$ signifie que l'on a $Z\ni h_i$ ou $Z\subset H_i$. \deleted{Il est} \inserted{Supposons le} invariant par tous les $\tau_i$, \deleted{ou bien $Z\subset H_i$ pour tt $0\leq i\leq n$, donc $Z\subset\bigcap H_i=\{c\}$, donc $Z=\varnothing$ ou $Z=\{c\}$, ou bien $\exists\,i$}

% L11–14
Notons que si $0<i\leq\unsure{n|n+1}$ $Z\subset H_i\Longrightarrow Z\subset H_{i-1}$, car sinon on aurait $h_{i-1}\in\unsure{Z|H_i}$ donc $h_{i-1}\notin H_i$, absurde. \unsure{De même|D'ailleurs}, si $h_j\in Z$, \deleted{\unsure{$\forall j$|$h_j$}} $0\leq j\unsure{<n|\leq n}$, alors $h_{j+1}\in Z$, car sinon on aurait $Z\subset H_{j+1}$ donc $h_j\notin H_{j+1}$, absurde. \deleted{\unsure{D}}Il en résulte \deleted{\unsure{qu'il}} \inserted{qu'on} a deux possibilités

% L17–18
a) $Z\subset H_i\qquad \forall\ 0\leq i\leq n$, donc $Z\subset\bigcap_{0\leq i\leq n}H_i=\{c\}$, donc $Z=\varnothing$, ou $Z=\{c\}$ -- qui sont bien invariants

% L19–21
a$'$) $Z\not\subset H_i\qquad \forall\ 0\leq i\leq n$ donc $Z\ni h_i\quad 0\leq\forall i\leq n$ donc $Z\supset\unsure{V|\vee}(h_0\,\ldots\,h_n)=C$, donc $Z=C$ ou \inserted{$Z=$}$X$ qui sont bien invariants

% L22–23
b) On n'est pas en cas a) ni b), donc $\exists\ i$, $0\leq i<n$, tel que $Z\subset H_j$ \inserted{si $0\leq j\leq i$} et $Z\not\subset H_j$ si $n\geq j\geq i+1$

% archive sheet number: 82
\scanpage{84}

% L1–5, upper right
Groupes engendrés par des pseudo-réflexions projectives et relations avec les réalisations géom. des \unsure{\textquotedblleft|\textquotedblright}\deleted{\unsure{cor}}hyperpolyèdres.

% archive sheet number: 83
\scanpage{85}

% L1, cancelled opening
\deleted{Soit $C$ $n$-maquette \unsure{rég.}\inserted{régulière}, \unsure{correspond|correspon}}

% L2–7
Soit $\unsure{C|\mathcal C}$ une $n$-carte régulière combinatoire, \deleted{\unsure{repérage}} $R$ l'ens. de \deleted{ses} repères, qui est un torseur à gauche sous $G=\operatorname{Aut}C$, le commutant de $G$ étant formé des opérations du groupe cartographique, engendré par $\sigma_0,\ldots,\sigma_n$. Si $\underline r\in R$, il lui correspond des générateurs $\unsure{\tau_i|\mathrm T_i}=\tau_i^{(\underline r)}$ de $G$, définis par

% L8
\[
\tau_i^{(\underline r)}\underline r=\sigma_i\underline r \tag*{(1)}
\]

% L9–11
Une réalisation géométrique projective de $C$, \unsure{dans|ds} $(n+1)$-fibré projectif $X$ sur $S$, revient à la donnée de a) Opération $\varphi$ de $C$ sur $X$

% L12
\[
\varphi:G\longrightarrow\operatorname{Aut}_S X=\Gamma\mathop{\underline{\operatorname{Aut}}}\nolimits_S(X)\tag*{(2)}
\]

% L13–14
b) drapeau \inserted{max} $r=(\deleted{\unsure{\Phi_0}}X_0,\deleted{\unsure{\Phi_1}}X_1,\ldots,X_{\unsure{n|n+1}})\deleted{X}$ dans $X$ \deleted{\unsure{avec}} satisfaisant

% L15–16
\[
\begin{aligned}
\tau_iX_j&=X_j\quad\text{si } i\ne j\\
\deleted{\tau_iX_i}&\deleted{\ne X_i}
\end{aligned}\tag*{(3)}
\]

% L17
(4) Les $\tau_i$ sont des réflexions de $P$

% L18
[ Ceci permet de \inserted{définir $R\longrightarrow\text{Drap.max}(X)$ et de} factoriser en $\operatorname{Drap}_{\unsure{+1|\ne1|H}}(C)\longrightarrow\operatorname{Drap}_{H}(X)$ ]

% L19
\uline{Prop} Conditions équivalentes (\deleted{\unsure{si on est}} \unsure{$S=\operatorname{Spec}(k)$|$S=\operatorname{Spec}(K)$} \inserted{$k$ corps})

% L20–22
b) Aucune des applications

% L21
\[
\deleted{\unsure{(4)}}\qquad \Phi_i\longrightarrow\operatorname{Grass}_i(X)\qquad 0\leq i\leq n
\]

% L22
n'est constante, \deleted{sur aucune fibre géom}

% L23
\deleted{b$'$) $\forall$ \unsure{$i$}, $0\leq i\leq n$, on a $\tau_i X_i\ne X_i$.}

% L24, left brace
\margin{indépendant du drapeau $r$.}
% slanted

% L24
d) si $\underline f_{i-1}<\underline f_i\ne\underline f'_i<\underline f_{i+1}$ ds $\Phi_*(C)$, alors $f_i\ne f'_i$

% L25–26
a) L'image de $\underline\Phi_0$ engendre lin$^{\mathrm t}$ $X$ (i.e. l'orbite de $X_0$ engendre lin$^{\mathrm t}$ $X$

% L26
\uline{a$'$)} L'intersection des hyperplans $\in\Phi_n$ est vide

% L27
\deleted{\unsure{d)}} c) $\forall\ 0\leq i\leq n$, $\tau_iX_i\ne X_i$

% archive sheet number: 84
\scanpage{86}

% L1–2, top right
\margin{\uline{NB} Dans f), les conditions $\alpha)$ $\beta)$ $\gamma)$ ne font intervenir que $(\unsure{\tau_i|\mathrm T_i})$, non le choix de $s_0$, qui est précisé dans $\delta)\,\unsure{\eta|\gamma})\,\unsure{\varphi|\xi|\rho})$.}

% L2–7
e) Posant

% L3–7
\[
\begin{aligned}
s_0&=\text{unique pt de }X_0\\
s_1&=\tau_0s_0\\
s_2&=\tau_1s_1\\
&\ldots\\
s_{n+1}&=\tau_n s_n,
\end{aligned}
\]

% L4–8, upper right
\margin{\uline{NB}\quad (5) $s_i=\deleted{\unsure{\tau_i}}\tau_{i-1}\ldots\tau_0s_0$ $(1\leq i\leq n+1)$\\
\textcircled{6} donc si $j>i$ $\tau_js_i=(\tau_{i-1}\ldots\tau_0)\tau_js_0=s_i$ car $\tau_js_0=s_0$.}

% L3–10, left diagonal
% slanted
\margin{il y a une condition duale e$'$), où $\deleted{X_n}S_{n+1}=X_n$, $S_{\unsure{n|n-1}}=\tau_n X_n$, $S_{n-1}=\tau_{n-\unsure{2|1}}S_{n-1}$ \ldots\ $S_{\deleted{\unsure{0}}1}=\tau_0 X_0$, en exigeant $\bigwedge_{0\leq i\leq n+1}S_i=\varnothing$}

% L8–9
les $s_i$ sont \deleted{\unsure{affin.t}} projectiv$^{\mathrm{nt}}$ indépendants, i.e. engendrent l'espace $X$.

% L10
[ f) \deleted{Soit} \deleted{\unsure{$h$}} $h_i$ ($H_i$) le centre (l'antre) de $\tau_i$. On a :

% L10–17, right
\margin{$\alpha)$ $\bigvee h_i$ hyperplan $C$\\
$\beta)$ $\bigwedge H_i$ pt $c$ ($H_i$\unsure{/|,|)} indép.)\\
$\gamma)$ $h_i\notin H_{i+1}$ $(0\leq i\leq n-1)$\\
$\delta)$ $s_0\in\underbrace{H_1\cap\cdots\cap H_n}_{(\text{droite }D)}$\\
(NB il est plus fort a priori que $\tau_i s_0=s_0\quad1\leq i\leq n$)\\
$\unsure{\eta|\gamma})$ $s_0\notin H_0$ ou encore $s_0\ne c$ ($\{c\}=\bigwedge H_i$) ou encore $\tau_0s_0\ne s_0$\\
$\unsure{\varphi|\xi|\rho})$ $s_0\notin C$ i.e. $s_0\ne d$ où $\{c'\}=D\cap C$.}

% L11
\uline{Schéma de dém} On a les implications

% L12–14, implication scheme
\diagram{\begin{tikzcd}[ampersand replacement=\&,column sep=small,row sep=small]
a) \arrow[dr,Rightarrow,"{\mathrm{tr}}"] \& \& \& \& a)\\
\& b) \arrow[r,Leftrightarrow,"{\mathrm{tr}}"] \& c) \arrow[r,Rightarrow] \arrow[d,Leftrightarrow,"{\mathrm{tr}}"'] \& e) \arrow[ur,Rightarrow,"{\mathrm{tr}}"] \arrow[dr,Rightarrow,"{\mathrm{tr}}"'] \&\\
a') \arrow[ur,Rightarrow,"{\mathrm{tr}}"] \& \& d) \& \& a')
\end{tikzcd}}

% L15
la seule implication non triviale étant $d)\Longrightarrow e)$

% L16–18
On prouve l'indép. des $(s_0,\ldots,s_i)$ par \deleted{\unsure{réc}} sur $i$, en \deleted{Si elle est établie} prouvant en \unsure{même} temps

% L19
\uline{Cor} Sous les conditions précédentes, on a

% L20–21
\[
\left\{\begin{aligned}
X_i&=\operatorname{Lin}(s_0,\ldots,s_i)\\
&\deleted{\tau_j s_i=s_i\ \text{si }j\unsure{>|\ne}i}
\end{aligned}\right.\quad
\begin{aligned}&(0\leq i\leq n+1,\deleted{\unsure{S}}\\[-2pt]
 &\text{posant }X_{n+1}=X)\end{aligned}
\tag*{(7)}
\]

% L22–23
Si c'est prouvé pour \deleted{$i$, \unsure{savoir}, $\tau_i$} $i\leq i_0$, prouvons le pour $i=i_0+1$, i.e.

% L24–26
\[
\left\{\begin{aligned}
X_{i+1}&=\operatorname{Lin}(s_0,\ldots,s_i,s_{i+1})\\
&\quad\bigl(\operatorname{Lin}(X_i,s_{i+1})\\
\tau_j s_{i+1}&=s_{i+1}\quad\text{si }j>i+1
\end{aligned}\right.
\]

% L26–27, extension to right
\margin{$\longrightarrow$ \deleted{\unsure{vérifié par l'hyp. de récurrence|vérifié par (5)}} \unsure{est|car} $s_i\in X_i\subset X_{i+1}$}

% L27–31
Par déf $s_{i+1}=\tau_i s_i$, or $\tau_i X_{i+1}=X_{i+1}$ donc \uline{$s_{i+1}\in X_{i+1}$}. \deleted{\unsure{On|Reste}} \inserted{Il faut prouver qu'on} ne peut avoir $s_{i+1}\in X_i$ i.e. $\tau_i s_i\in X_i$, car \deleted{si} $X_i=\operatorname{lin}(X_{i-1},s_i)$ donc $\tau_iX_i=\operatorname{lin}(\underbrace{\tau_i X_{i-1}}_{X_{i-1}},\underbrace{\tau_i s_i}_{\unsure{\text{OK}}})$ qui serait égal à $X_i$, contraire à l'hyp. \unsure{OK|c)}

% archive sheet number: 85
\scanpage{87}

% L1–2, left
\[
\begin{aligned}
\unsure{a_0|a.}\quad\ldots\quad a_{i-1} && a_{n-2}\ \overline{a_{n-1}\ a_n}\\
&&a'_{i+1}\quad\ldots a'_{n-1}\ a'_n
\end{aligned}
\]

% L1–4, upper right
\[
\begin{aligned}
c&=\bigwedge H_i & d&=\bigl(\overbrace{\bigwedge_{1\leq i\leq n} H_i}^{Z_1}\bigr)\wedge C\\
C&=\bigvee h_i & D&=\underbrace{\bigvee_{0\leq i\leq n-1}h_i}_{Y_{n-1}}\vee c
\end{aligned}
\]

% L3–4, left
\[
\begin{aligned}
Y_{n-1}\cap H_n&=Y_{n-2}\\[-2pt]
&\ldots\\
Y_{n-1}\cap(H_n\cap\ldots\cap H_{i+1})&=Y_{i-1}
\end{aligned}
\]

% L5–6, right box
\[
\boxed{\begin{aligned}Z_1&\not\subset C\\ c&\notin Y_{n-1}\end{aligned}}
\]

% L5–11, central ladder
\diagram{\begin{tikzcd}[ampersand replacement=\&,row sep=small,column sep=small]
Y_{n-1}\arrow[d,dash,"{n-i}"'] \& \overset{\overline{\scriptstyle 2}}{\subset} \& X\arrow[d,dash,"{n-i}"]\\
Y_{i-1}\arrow[d,dash,"{n}"'] \& \underset{2}{\subset} \& Z_{i+1}\arrow[d,dash,"i"] \& {}=H_n\cap\ldots\cap H_{i+1}\\
\varnothing \& \underset{2}{\subset} \& Z_1
\end{tikzcd}}

% L7–9, small diagram
\diagram{\begin{tikzcd}[ampersand replacement=\&,column sep=small,row sep=small]
\deleted{c,d} \& Y_0\arrow[d,dash] \& h_0\notin H_1\\
\&\varnothing\arrow[r,dash]\&\cdot\\
\&c\quad d\&
\end{tikzcd}}

% L9–11, right box
\[
\boxed{\begin{aligned}
X_n&\longmapsto X_i=X_n\wedge Z_{i+1}\\
X_0&\longmapsto X_i=X_0\vee Y_{i-1}
\end{aligned}}
\]

% L12–13
\[
\begin{aligned}
c&=\bigwedge_0^n H_i &d&=(\overbrace{\bigwedge_1^n H_i}^{Z_1})\wedge C=Z_1\wedge C\\
C&=\bigvee_0^n h_i &D&=\bigl(\underbrace{\bigvee_0^{n-1}h_i}_{Y_{n-1}}\bigr)\vee c=Y_{n-1}\vee c
\end{aligned}
\]

% L14–16, left
\[
\begin{aligned}
X'_i&=\deleted{\unsure{\bigl(}}(Y_{i-1}\vee c)- & X''_i&\unsure{=|:=|\overset{?}{=}}\deleted{\unsure{\{} }Y_{i-1}\vee d\\
&\deleted{\unsure{C\wedge Z_{i+1}|=C\wedge Z_{i+1}}}& &= Z_{i+1}\wedge D\\
&=Z_{i+1}\wedge C
\end{aligned}
\]

% L14–16, right
\[
\begin{aligned}
Y_{n-1}&\ \deleted{\unsure{\sim|-}}\ Z_1\\
Y_1&\qquad Z_{n-1}=H_n\wedge H_{n-1}
\end{aligned}
\]

% L17–20, left
\diagram{$\begin{array}{cc} h_i&H_i\\ c&C\\ d&D\end{array}$}

% L17–21, right box
\[
\boxed{\left\{\begin{aligned}
 h_i&\subset H_j &&\text{si }i,j\ \text{non consécutifs}\\
 h_i&\notin H_{i+1}\\
 &\deleted{(h_i)\ \text{indép.}}\\
 \bigvee h_i&=C &&\text{de codim }1\\
 \bigwedge H_i&=c &&\text{de dim }0\quad /\,d
\end{aligned}\right.}
\]

% L21, far left
\deleted{$\unsure{\bigwedge}h_i$}

% L22–23, bottom centre
\diagram{$\begin{array}{cc}s_0&d\\\text{l'origine}&\text{infini}\end{array}$}

% archive sheet number: 86
\scanpage{88}

% L1–3
\deleted{\unsure{$h_0\,h_1$}} $h_0\ h_1\ \ldots\ h_n$ sections du fibré projectif $X$. $H_0\ H_1\quad H_n$ \deleted{hyperplans} \deleted{\unsure{$X$}}

% L1–6, right
\margin{$(h_i,H_i)$ centre et cocentre \deleted{de} de pseudo réflexion $\unsure{\tau_i|\mathrm T_i}$ dans $X$\\ $\tau_i\tau_j=\tau_j\tau_i$ si $i,j$ non cons.\\ d'où $\unsure{h_i|H_i}$}

% L3–4
\[
\begin{aligned}
\varnothing=X_{-1}\subset X_0=\{s_0\}&\subset X_1\subset\ldots\subset X_n\subset X=X_{n+1}
 & & (X\text{ de rang }n+1)\\
&&&\tau_i X_j=X_j\ \text{si }i\ne j\quad 0\leq i,j\leq n\\
&\Longrightarrow &&\underline\Phi_i\longrightarrow\Phi_i
\end{aligned}
\]

% L5–7
A) Conditions équivalentes

% L6–7, crossed opening
\deleted{a) $\forall\ 0\leq i\leq n$, $\Phi_i$ ne se réduit à un pt b) $\forall\ 0\leq i\leq n$}

% L7–8
a) $\Phi_0$ engendre lin. $X$ en chaque fibre

% L8
a$'$) l'intersection des $Z\in\Phi_n$ est vide (\deleted{\unsure{ds}} \inserted{en} chaque fibre)

% L9
b) $\forall\ 0\leq i\leq n$, $\Phi_i$ non \unsure{réduite|réduit} à un pt

% L10
c) $\forall\ 0\leq i\leq n$ $\unsure{\tau_i|\Phi_i}(X_i)\ne X_i$\qquad $\exists$

% L11–12
d) si $\underline f_{i-1}<\deleted{\underline f_i\ne\underline f'_i}<\underline f_{i+1}$ dans $\underline\Phi_*$, alors $\underline f_i\ne\underline f'_i$ entre $\underline f_{i-1},\underline f_{i+1}$ tels que $f_i\ne f'_i$.

% L13–14
e) Posant $s_1=\tau_0 s_0$, $s_2=\tau_1s_1=\tau_1\tau_0s_0$, \ldots, $s_{n+1}=\tau_ns_n=\tau_n\tau_{n-1}\ldots\tau_0s_0$ les $(s_i)_{0\leq i\leq n+1}$ engendrent \deleted{affinement} \inserted{lin$^{\mathrm t}$ $X$}

% L15–16
e$'$) Condition duale sur $S_n\inserted{=\text{ par déf }X_n}$, $\deleted{\unsure{\tau_n S_n=}}S_{n-1}=\tau_n S_n$, $S_{n-1}=\tau_{n-1}S_{n-1}=\tau_{n-1}\tau_n S_n\ldots$ $S_{-1}=\tau_0S_0=\tau_0\tau_1\ldots\tau_n S_n$, savoir $\bigwedge_i S_i=\varnothing$.

% L17
\deleted{f) L'une des conditions} f) On a l'ens des \deleted{\unsure{ces}} conditions

% L18–20
\[
\left\{\begin{aligned}
\alpha)&\quad \deleted{\unsure{\text{les}}}(h_i)_{0\leq i\leq n}\ \text{libres \inserted{en tte fibre} -- engendre hyperplan }C\\
\alpha')&\quad(H_i)_{0\leq i\leq n}\ \text{libres \inserted{en ttes fibres} -- }\bigcap_j H_i=\{c\},\ c\ \text{section de }X\\
\beta)&\quad h_i\unsure{\notin|\not\subset}H_{i+1}\ \text{en chaque fibre}\quad 0\leq i\leq n-1
\end{aligned}\right.
\]

% L18–25, left box
\margin{\uline{NB} $\alpha)\,\alpha')$ impliquent $\left\{\begin{smallmatrix}h_i\subset H_j\\h_j\subset H_i\end{smallmatrix}\right.$ si $0\leq i,j\leq n$, $i,j$ non consécutifs (ce qui exprime la cond. de commutation des $\tau_i$) $f)$ implique $X_i=\deleted{\unsure{V|\bigvee}}V(s_0,h_0,\ldots,h_{i-1})=\bigwedge(X_n,H_n,\ldots,H_{i+1})$ $X_i$ et $X_j$ se déterminent mutuellement pour $i<j$ par $X_i=\bigwedge(X_j\cap H_j\cap\ldots\cap H_{i+1})$, $X_j=V(X_i,h_i,\ldots,h_{j-1})$}

% L21–25, right, overwritten
\deleted{\unsure{$\gamma)$}} posons \deleted{\unsure{$D$|$D_0$}} $D_0=H_n\cap\ldots\cap H_1$ (droite), on a déjà a) \deleted{$D_0\cap C=\{d\}$, \unsure{$c$ et $d$|$c\in D_0$}} \inserted{$D_0\cap C=\deleted{\unsure{\{d\}}}\ \text{une section }\{d\}$} \deleted{\illegible{2} \unsure{fixe et $d$} \illegible{2} \unsure{de $\tau_0$})}

% L23–25
$\delta)$ $s_0$ section de \deleted{\unsure{$D_0$}} $D_0^{**}=D_0\smallsetminus\{c\cup d\}=D_0-$\deleted{sections de $X$} pts fixes de $\tau_0\mid D_0$

% L26–28
B) Pour $(\tau_i)$ donné, pour qu'il existe loc. (Zar) un drapeau $(X_i)$ satisfaisant aux conditions équivalentes ci-dessus, il f et s. que $(\tau_i)$ satisfasse les conditions $(\alpha,\alpha',\beta)$ de (A f).

% L29–32
C) Sous les conditions de B (i.e. $\alpha,\alpha',\beta$), l'application $(X_i)\longmapsto s_0$ est une bijection de l'ens \inserted{$\operatorname{Dr}(\tau_*)$} des drapeaux $(X_i)$ adaptés à $(\tau_i)$ et non dégénérés, avec l'ens des sections de $\unsure{\delta_0^{**}|D_0^{**}}$. Plus gén., \deleted{\unsure{$i$ donné}}, \inserted{$\forall\ 0\leq i_0\leq n$},

% L32–36
\deleted{bijection de $\operatorname{Dr}(\tau_*)$ sur l'ens} $(X_{\unsure{j|\nu|*}})\longmapsto X_{i_0}$ est une bijection de $\operatorname{Dr}(\tau_*)$ sur l'ens des sections de $D_i^{**}$, où $D_i^{**}=D_i\smallsetminus\{c_i\cup d_i\}$, $D_i=$ fibré en droites des $(i+1)$-espaces lin. compris entre $Y_{i-1}=V(h_0,\ldots,h_{i-1})$ et $Z_{i+1}=\bigwedge(\deleted{\unsure{X|H_n}}H_n,\deleted{\unsure{\Phi}}\ldots,H_{i+1})$ (resp. de dim $i-1$ et $i+1$), $c_i$, $d_i$ étant les sections de $D_i$ définies resp. par

% L37–39, bottom crossed block
\deleted{$Y''_i=\unsure{Y_{i-1}|Y_{i-1}\vee Z_{i+1}}$, $Z_{i+1}$ $Y'_i$ et $Z'_i=\bigvee(Z_i,Y_{i-1})$ $Y_{i-1}\subset Y'_i,Z'_i\subset Z_{i+1}$ [\uline{NB} $Y_i$ stable par $\tau_0\ldots\tau_i$, $Z_i$ stable par $\tau_i,\ldots,\tau_n$].}

% L38–40, bottom interline
\[
\begin{aligned}
X'_i&=\deleted{\unsure{V}}(Y_{i-1}\vee c) &\text{et}\quad &\deleted{\unsure{V|\bigl(}}Y_{i-1}\vee d\deleted{\unsure{\}|\bigr)}}=\overset{\deleted{\unsure{\text{c'est aussi}}}}{X''_i}\\
&\ \Vert &&\ \Vert\\
&\deleted{\unsure{X}}(Z_{i+1}\wedge \unsure{D|C}) && \deleted{\unsure{X}}(Z_{i+1}\wedge D)
\end{aligned}
\]

% L40, bottom right
\margin{où $D$ est défini dualem$^{\mathrm t}$ à $c$ comme $D=V(h_1\ldots h_{n-1},c)$}

% archive sheet number: 87
\scanpage{89}

% L1–5
D) Sous les conditions de B), soit \deleted{\unsure{$G$ de}} $\underline g$ le groupe des autom de $(X,\deleted{(\unsure{\tau_*|\tau_i})})$, i.e. des $G_n$-autom de $X$, où $G_n$ est le groupe engendré par les gén. $\underline\sigma_i$ ($0\leq i\leq n$), liés par relations $\underline\sigma_i\underline\sigma_j=\underline\sigma_j\underline\sigma_i$ si $i,j$ non consécutifs. On a \deleted{donc}

% L6
\[
\underline g\longrightarrow\underline{\operatorname{Aut}}\bigl(D_0,\overset{\text{id sur}}{(c\sqcup(d))}\bigr)
=\operatorname{Aut}\bigl(D_0,\overset{\text{id sur}}{(\Delta_0)}\bigr)
\]

% L5–7, interline right
\margin{produit de div. relatifs}

% L6–8, right
\margin{où $\Delta_0$ est le diviseur sur $D_0$ défini par les sections $c,d$}

% L6–10, left
\margin{\uline{NB} $\underline g$ est lisse \unsure{comm.|connexe|commutatif} rel. \unsure{conn.|un.} de dim rel $1$, ses fibres sont isomorphes à $\mathbb G_m$ ($c\ne d$) ou $\mathbb G_a$ ($c=d$).}

% L7–12
Ceci dit, cet hom est un \uline{isom} -- par suite $D_0^{**}\inserted{=D_0\smallsetminus\Delta_0}$ est un $\underline g$-torseur : $\exists$ loc sections (c'est le contenu de B,C), et deux sections sont conjuguées par un unique \inserted{$G_n$-}autom. Donc \deleted{\unsure{en|En}} en vertu de C, itou pour le schém. des \inserted{$(\tau_*)$-}drapeaux non dégénérés, qui est un torseur sous $\underline g$.

% L13–19
\uline{Scholie} Un système $(\tau_*,X_*)$ \inserted{non dégénéré} peut être considéré comme une \uline{rigidification} d'un système $(\tau_*)$ non dég, i.e. d'une représentation projective n. dég. de $G_n$. La cat. (sur $S$ fixé) des $(\tau_*\deleted{\unsure{,X_*}})$ non dég. est équivalente à celle des $(\tau_*,X_*)$ (i.e. des \uline{réal. géom non dég de $G_n$}) munis d'un torseur sous le groupe \inserted{$\unsure{g_0|\underline g_0}$} $\underline{\operatorname{Aut}}(D_0,\overset{\text{id sur}}{\unsure{c_0;d_0|\Delta_0}})$ correspondant à cette dernière\ldots\ Utilisant la section $s_0$ disjointe de $d$, \deleted{comme \unsure{et} $d$ comme} \inserted{le normant pt $1$, $c$ comme $0$ et $d$ comme $\infty$}

% L20–24
(si $c\ne d$ en tte fibre i.e. le centre n'est à l'$\infty$ sur aucune fibre) on trouve que ce groupe est $G_m$, donc \deleted{il s'agit} \inserted{dans ce} \deleted{de} \inserted{cas} \deleted{\unsure{don|la}} la cat. des $(\tau_*)$ non dég \unsure{non|son|avec} centre partout à dist. finie équivaut à la cat produit de la cat des $(\tau_*,X_*)$ par celle des Mod. inv. sur $S$.

% L25–26
E) Soit $\Pi$ un quasi-polyèdre combinatoire régulier. $G$ le groupe de ses autom, $R$ l'ens. des repères, qui

% archive sheet number: 88
\scanpage{90}

% L1–4
est un torseur sous $G$. Le choix d'un $r\in R$ définit un \deleted{\unsure{épim|épim.}} $G_n\xrightarrow{\varphi_r}G$, d'où un système \inserted{$\unsure{\tau_*|\mathrm T_*}$} $\underline\tau_*^r=(\underline\tau_i^r)=\varphi_r(\underline\sigma_i))$, transformé en un syst. conjugué \inserted{par $\operatorname{int}(g)$} si on change $r$ en $r'=gr$ ($g\in G$). \unsure{Si|\deleted{Si}} on se donne un hom

% L5
\[
\varphi_G:G\longrightarrow\operatorname{Aut}X,
\]

% L6–8
dire que $\tau_*^r=(\tau_i^r)=\varphi_G(\underline\tau_i^r)=(\varphi_G(\underline\tau_i^r))$ est formé de \uline{quasi-réflexions}, et que la condition de non dég. de B est satisfaite, ne dépend pas du choix de $r$.

% L9–11
De $=$ si on a une repr. \unsure{régulière} $(\varphi_G,\inserted{\varphi_*=}(\varphi_i)_{0\leq i\leq n}\quad(\varphi_i:\Phi_i(\Pi)\longrightarrow\operatorname{Grass}_i(X)))$ de $\Pi$, d'où $r=\varphi_*(\underline r)=X_*$, dire que $(\tau_{\unsure{j|*}}^r,X_*))$ \deleted{est} est non dég ne dépend pas du choix de $\underline r$ -- \deleted{satisfaire} \inserted{d'ailleurs}

% L12–16
les conditions \inserted{a) a$'$) b)} d) de A) sont intrinsèques. La première est \unsure{donc|en vertu} néc. et suff. pour l'existence loc$^{\mathrm{le}}$ (Zar) d'une \textquotedblleft réal. géom non dég\textquotedblright\ de $\Pi$ associée à $\varphi_G$. \deleted{\unsure{De même,|Donc} Ici encore, on voit que le groupe} \inserted{Comme $G$ est \uline{simplem$^{\mathrm t}$} \uline{transitif} sur} $R$, on voit d'ailleurs que tt objet géom. qu'on

% L17–20
peut associer intrinsèquement à une situation $\tau_*$ resp $(\tau_*,X_*)$ non dég [ tel que $D_0$, et ses sections $c,d$ ] est défini intrinsèquement aussi en termes d'une \inserted{repr.} $\varphi_G$ \inserted{quasi-polyédrale n.d.} resp. d'une \deleted{\unsure{$\varphi$}} réalisation géom \inserted{$(\varphi_G,\varphi_*)$} de $\Pi$.

% L21–22
Ceci dit, $\operatorname{Aut}\varphi_G$ opère sur $(D_0,c,d)$, et on trouve

% L22
\[
\underline{\operatorname{Aut}}(X,\varphi_G)\xrightarrow{\sim}
\underline{\operatorname{Aut}}\bigl(D_0;\overset{\text{id sur}}{(c)\cdot(d)}\bigr)
\]

% L23–25
donc la cat. des $\varphi_G$ proj.le non dég. équivaut à celle des couples $\deleted{\unsure{(\varphi_G}}\bigl(\underbrace{(\varphi_G,\varphi_*)}_{\substack{\text{réal. géom proj.}\\\text{n. dég. de }\Pi}},T\bigr)$, $T$ un torseur sous $\underline{\operatorname{Aut}}(D_0,c,d)$ -- ou la \unsure{$=$|même} interprétation

% archive sheet number: 89
\scanpage{91}

% L1–7, upper quarter
que ci-dessus pour $T$ comme l'ens des structures \unsure{qu. pol.|quasi-pol.} non dég. subordonnées à $\varphi_G$. Lorsque $c\ne d$ en chaque pt, $T$ s'interprète comme un torseur sous $G_m$ (opérant par homothéties, cf plus bas) Si $c=d$, alors $D_0^*\inserted{=D_0\smallsetminus c}$ est un \deleted{\unsure{tor|torseur} fibré affine}, \deleted{trivialisé par $s_0$} \deleted{un fibré vectoriel} de rang $1$, \inserted{trivialisé par $s_0$} et $\unsure{\simeq D_0^*|\simeq D_0^\times}$ le fibré vectoriel \unsure{ass.|inv.} de ses translations\ldots

% archive sheet number: 90
\scanpage{92}

% L1
\uline{Coordonnées universelles} \inserted{$(n\geqslant1)$}

% L1–4
On prend. $C=V(h_0,\ldots h_n)$ comme hyperplan à l'$\infty$, \deleted{alors les $s_i$ \unsure{$\in E$|$(i\leq)$}} donc $E=X\smallsetminus C$ comme fibré affine de rg $n+1$, les $s_i$ sont donc des sections de $E$ formant une base affine

% L5–7
\deleted{et de \unsure{même} les $s_i$\unsure{$-h_0$|$-s_0$} \unsure{les} forment une base du fibré vectoriel des transl. de $E$ en termes de choix. On choisit comme base $s_0,s_1\ldots s_n$}

% L8–10
Le vectoriel $\unsure{\mathcal E|E}$ ext. de $\mathcal O_S$ par \deleted{\unsure{$\mathcal O_S$}} l'esp. des tr. $V$ de $E$ \deleted{\unsure{par une}} \inserted{enveloppe} vect. de $E$ s'identifie donc à $\mathcal O_S^{n+2}$, au vecteur \deleted{$(\lambda_0,\lambda_1,\ldots,\lambda_{n+1},\unsure{\lambda_{n+2}})$ correspond} \inserted{$(\lambda_0=1,\lambda_1,\ldots,\lambda_{n+1})$ correspond}

% L11
\[
\deleted{s_0+\lambda_1(\underbrace{s_1-s_0}_{e_0})+\lambda_2(\underbrace{s_2-s_0}_{e_1})+\ldots+\lambda_{n}(\underbrace{s_n-s_0}_{e_n})}
\]

% L12–13
\[
\begin{aligned}
\mathcal E=k s_0\oplus k s_1\oplus\ldots\oplus k s_{n+1}&\xrightarrow{\unsure{\pi|\Phi}}\mathcal O_S\\
\sum\lambda_i s_i&\longmapsto\sum\lambda_i
\end{aligned}
\]

% L14–15
$\unsure{\tau_i|\mathrm T_i}$ \inserted{\unsure{ps.|pso.|ps.-}}réflexions de $\mathcal E$ \inserted{$0\leq i\leq n$} satisfaisant

% L16
\[
\deleted{\tau_i\tau_j=\tau_j\tau_i\ \text{si }i\ne j}\qquad \tau_i=\mathrm{id}+a'_i\otimes a_i
\]

% L17–20, bracketed conditions
\[
\begin{aligned}
\deleted{\tau_i s_i=}\qquad\text{\textcircled{1}}\quad&
 \langle\deleted{\unsure{a}}s_i,a'_j\rangle=0,\quad\text{si }\deleted{i,j\ \text{non cons.}}\ \inserted{0\leq i<j\leq n},\ \deleted{j\ne i+1}\\
\text{\textcircled{3}}\quad&
 \deleted{\langle s_0,a'_j\rangle=0\quad\text{si }j\geq1}\ 
 \inserted{\langle a_j,a'_i\rangle=0\quad\text{si }0\leq i\unsure{<j|,j|\leq j}\leq n,\ j\ne i+1}\\
\deleted{\text{\textcircled{3}}\quad}&
\deleted{\langle s_0,a'_0\rangle,\langle a_0,a'_1\rangle,\ldots,\langle a_{n-1},a'_n\rangle\ \text{inversibles}}\\
\text{\textcircled{2}}\quad& s_{i+1}=\tau_i s_i\quad0\leq i\leq n\deleted{\unsure{+1}}
\end{aligned}
\]

% L18, left
% slanted
\margin{moins importante $\longleftrightarrow$}

% L21
\[
\underline{a'_j=\lambda_{jj}s'_j+\lambda_{j,j+1}s'_{j+1}\ldots+\lambda_{j,n+1}s'_{\unsure{j+1|n+1}}}\qquad\text{exprime (1) :}
\]

% L22–23
\[
\begin{aligned}
(\tau_i s_i=)\ s_i+\lambda_{ii}a_i&=s_{i+1}\qquad\deleted{a_i=}\lambda_{ii}a_i=s_{i+1}-s_i\\
&\text{ops }\lambda_{ii}=1,\ \text{donc}
\end{aligned}
\]

% L24–25
\[
\begin{aligned}
\boxed{a_i=s_{i+1}-s_i}\qquad\lambda_{ii}=1\quad &\text{i.e.}\quad a'_i=\unsure{s'_i|\beta'_i|s_i}+\lambda_{i,i+1}s'_{i+1}\ldots+\lambda_{i,n+1}s'_{n+1}\\
&\text{expriment (1) et (2)}
\end{aligned}
\]

% L26–27
\[
\left.\begin{aligned}
\lambda_{i,i+2}=\lambda_{i,i+3}=\ldots&=\lambda_{i,\deleted{\unsure{n+}}n+1}
 \quad(\unsure{\text{soit}}\ 1+\unsure{\mu_i|\alpha_i})\\
\lambda_{i,i+1}&=\lambda_i
\end{aligned}\right|\quad\text{exprime \uline{(3)}}
\]

% archive sheet number: 91
\scanpage{93}

% L1–2
on peut supposer ces conditions satisfaites, \deleted{\unsure{donc}} \inserted{donc aussi (9)} \inserted{donc aussi (10) vrai.} \deleted{On \unsure{aura}} \inserted{Notons} qu'alors \deleted{l'indépendance de} la condition

% L3–6
$s_{n+1}$ \deleted{\unsure{$f$|$\notin$}} libre mod $F_n$ équivaut à la condition \deleted{libre mod $F_n$} (5), \deleted{Maintenant on sait déjà que $a'_n\mid F_{n-1}=0$, donc} \inserted{Mais on sait qu'elle} équivaut à $\langle s_n,a'_n\rangle$ \deleted{$\in$ \unsure{inversible|inv.}}, et $\langle s_n,a'_n\rangle=\lambda_n$, OK.

% L7–9, diagonal cancellation begins
\deleted{Il reste, pour compléter la récurrence, à établir (8) pour $i=n+1$, or on \unsure{aura|a}}

% L10, cancelled display
\[
\deleted{s_{n+1}\overset{\text{déf}}{=}\tau_n s_n=s_n+\langle s_n,a'_n\rangle a_n\equiv\langle s_n,a'_n\rangle a_n\qquad(F_n)}
\]

% L11–14
\deleted{et comme $\underset{\in F_n}{s_n}$ \unsure{$s_n\equiv$} $\lambda_{n-1}a_{n-1}\langle s_{n-1},a'_n\rangle a_n\equiv\langle s_n,a'_n\rangle a_{n-1}$ mod $F_{n-1}\subset F_n$ donc et \unsure{fait|mod} $F_n=\mathcal O_S(s_0,a_0,\ldots,a_{n-1})$ mod $F_n$, \unsure{on aura} \unsure{déjà acquis}}

% L13–16, nested workings
\[
\begin{aligned}
&\deleted{\left\{\begin{aligned}
s_n&\equiv\deleted{\unsure{s_{n-1}+}}\lambda_{n-1}a_{n-1}\qquad(F_{n-1})\\
a'_n&\mid F_{n-1}=0
\end{aligned}\right.}\\
&\deleted{s_{n+1}\equiv\langle s_n,a'}\\
&\deleted{s_{n+1}\equiv\lambda_{n+1}\lambda_{n-1}a_{n-1}}\qquad
\deleted{\overset{\Vert}{\lambda_n a_n}\quad\text{par (9)}}\\
&\deleted{\text{d'où }\langle s_n,a'_n\rangle=\unsure{\lambda_n.}}\qquad\deleted{\text{OK.}}
\end{aligned}
\]

% L15–17, left
% slanted
\margin{L'assertion de stabilité évidente\ldots}

% L18–19
\deleted{\uline{Corollaire} Pour fibré vectoriel ($\mathcal O$ Mod. l.l. de \unsure{r.f.|t.f.}) et conditions fibre par fibre.}

% L20–22
\deleted{\uline{Reformulation} \deleted{\unsure{aussi}} Cas où $E$ loc. libre de rang $n+2$. Conditions équivalentes sur le système $(\overset{\text{section}}{s_0},\overbrace{\unsure{\tau_0|\mathrm T_0},\ldots,\tau_n}^{\text{ps.réflexions}})$, définissant les $s_i$}

% L23–24, bottom upside down
% margin note written sideways (upside down)
\margin{$a'_i\otimes \underline{a_i}$\\ $\underline{a'_{i+1}}\otimes a_{i+1}$}

% L23–24, bottom upside down, right of products
% upside down
\diagram{\begin{tikzcd}[ampersand replacement=\&,row sep=small,column sep=small]
\deleted{\unsure{\xi}} \& a_0 \arrow[r,dash]\arrow[dr,dash,"{a'_1}"] \& a_n\\
\&\&{}
\end{tikzcd}}

% archive sheet number: 92
\scanpage{94}

% L1–6, slashed block
\[
\begin{aligned}
&\deleted{\underbrace{s_0,\unsure{a_0},\ldots,a_{i-1}}}\\
&\deleted{\langle s_0,a'_0\rangle=1}\\
&\deleted{\langle\unsure{a_i|h_i|s_i},a'_{i+1}\rangle=\langle s_{i+1}-s_i,s'_{i+1}+\ \rangle=1}\\
&\deleted{\langle a_j,a'_i\rangle=\langle s_{j+1}-s_j,s'_i+\lambda_{i,i+1}s'_{i+1}\ldots+\lambda_{i,n+1}s'_{n+1}\rangle}\\
&\deleted{\langle a'_i,a_j\rangle=}\qquad\deleted{j\geq i+2}\\
&\deleted{i\leq j-2}\qquad\deleted{\lambda_{i,}}\\
&\deleted{i+2\leq j}
\end{aligned}
\]

% L7–10, box
\[
\boxed{\begin{aligned}
\unsure{\tau_i|\mathrm T_i}&=\mathrm{id}+a'_i\otimes a_i\\
a_i&=s_{i+1}-s_i\\
a'_i&=s'_i+\lambda_i s'_{i+1}+(1+\alpha_i)(s'_{i+2}+\ldots+s'_{n+1})\quad 0\leq i\leq n-1\\
a'_n&=s'_n+\lambda_n s'_{n+1}
\end{aligned}}
\]

% L7–9, right diagonal
% slanted
\margin{paramètres $\overbrace{\lambda_0\ldots\lambda_n}^{n+1}$ $\underbrace{\alpha_0\ldots\alpha_{n-1}}_n$}

% L11
Quand a-t-on

% L11–12, right enclosure
\[
\begin{aligned}
Y_i&\subset Z_{i+1}\\
\overline{V(h_0\ldots h_{i-1})}&\qquad\bigwedge(H_n,\ldots H_{i+1})
\end{aligned}
\]

% L13–14
\[
\tau_i^2=1\qquad\text{i.e.}\qquad
\underset{\substack{\Vert\\\lambda_i-1}}{\langle a_i,a'_i\rangle}=-2\quad ?\qquad
\text{ssi }\boxed{\lambda_i=\unsure{-1|+1}}
\]

% L14–16
\[
\begin{aligned}
\tau_i&=\mathrm{id}+a'_i\otimes a^{\unsure{\prime|\deleted{\prime}}}_i\\
&\deleted{\unsure{a_0=}}\\
\tau_i^n&=\mathrm{id}+P_n(\langle a_i,a'_i\rangle)
\end{aligned}
\]

% L17
\[
\tau_i\tau_{i+1}=\mathrm{id}+a'_i\otimes a_i+a'_{i+1}\otimes a_{i+1}
 +\underbrace{\langle a_{i+1},a'_i\rangle}_{\alpha_i-\lambda_i\deleted{\unsure{+0}}}\,a'_{i+1}\otimes a_i
\]

% L18–19
\[
\begin{aligned}
\deleted{\unsure{\bigl( V_i}}&\qquad V_i\quad a_i,a_{i+1}\\
\deleted{\unsure{V_i\mid|V_i^\vee}}\longleftarrow &V'_i\quad a'_i,a'_{i+1}
\end{aligned}
\]

% L18–23, inverted block on right
% margin note written sideways (upside down)
\diagram{$\begin{array}{ccc}
a_0&\deleted{\unsure{a_1}}&a_1\\
a'_0&\deleted{\unsure{a'_1}}&a'_1\\
\deleted{s_0}&\deleted{s_1}&\deleted{\unsure{s_1\ne s_0}}\\
\multicolumn{3}{c}{\deleted{\tau_0s_0\ne s_0}}\\
\deleted{\unsure{a_{-1}}}&\deleted{a_0}&\deleted{\ldots a_n}\\
\deleted{a'_{-1}}&\deleted{a'_0}&\deleted{\ldots a'_n}
\end{array}$}

% L20–22, lower left
\[
\begin{aligned}
&\deleted{\underbrace{\lambda_i\lambda_{i+1}}+\lambda_i-\alpha_i}\\
&\deleted{4-1\,\alpha_i+\lambda_i}\\
&\underbrace{\unsure{+1|t}\begin{pmatrix}-2&1+\alpha_i\\1&-2\end{pmatrix}}_{\unsure{3|2}-\alpha_i}
\end{aligned}
\]

% L20–23, two matrices
\[
\begin{aligned}
&\begin{pmatrix}\lambda_i\deleted{\unsure{\alpha_{i+1}}}-1&1+\alpha_i-\lambda_i\\1&\lambda_{i+1}\deleted{\unsure{-\alpha}}-1\end{pmatrix}\\
&\begin{pmatrix}\langle a_i,a'_i\rangle&\langle a_{i+1},a'_i\rangle\\
\langle a_i,a'_{i+1}\rangle&\langle a_{i+1},a'_{i+1}\rangle\end{pmatrix}
\end{aligned}
\]

% archive sheet number: 93
\scanpage{95}

% L1
\uline{Lemme} Soit $\unsure{\mathcal E|E}$ fibré vectoriel de rg $N\geq3$. $X=\check{\mathbb P}(\mathcal E)$.

% L2–3
a) Soit $\unsure{\tau_X|\mathrm T_X}$ une \inserted{ps.}réflexion de $X$. Alors $\exists1$ \inserted{ps.}réflexion $\tau_{\mathcal E}$ de $\mathcal E$ qui lui donne naissance, propre \unsure{ssi|\deleted{ssi}|si} $\tau_{\deleted{\unsure{X}}\mathcal E}$ l'est.

% L4–7
b) Soit $C$ un \deleted{hyperplan} \inserted{sous-fibré proj} de $X$, \deleted{stable pour $\tau_X$, de} \inserted{$C=\check{\mathbb P}(V)$, supposons} \deleted{$\tau_X C=C$ i.e. $\tau_{\mathcal E}V=V$.} $\tau_X$ propre, centre $h\in C$, \deleted{alors} alors \deleted{Considérons l'extension $\mathcal E$ de $\mathcal O_S$ par $V$} $\tau_{\mathcal E}$ induit l'identité sur $\mathcal E/V$.

% L7–11
Si $C$ est hyperplan et $\mathcal E$ l'enveloppe vectorielle can. de $E=X\smallsetminus C$, donc $\mathcal E/V\simeq\mathcal O_S$ can$^{\mathrm t}$. alors l'action de $\tau_X\mid E$ \deleted{sur $E$} s'identifie à $\tau_{\mathcal E}\mid E$, si on identifie $E$ à l'image inv. de $1$ par $\mathcal E\longrightarrow\mathcal O_S$

% L12–13
Dégager : condition de stabilité d'un ss-fibré projectif par pseudo réflexion

% L14–15
conditions de commutation de deux pseudo-réflexions.

% L16
Ordre d'une pseudo-réflexion

% L17
Ordre du produit de deux pseudo-réflexions

% archive sheet number: 94
\scanpage{96}

% L1–3, upper right
Groupes engendrés par des pseudo-réflexions vectorielles\ldots

% archive sheet number: 95
\scanpage{97}

% L1–2
\uline{Théorème 1} Soient $k$ un anneau, $\unsure{\mathcal E|E}$ un $k$-module, $n\in\mathbb N$ \deleted{\unsure{et des|soient}} $\unsure{\tau_i|\mathrm T_i}$ $(0\leq i\leq n)$ des pseudo réflexions dans $\mathcal E$

% L3
\[\tau_i=\mathrm{id}+a'_i\otimes a_i\qquad a'_i\in\mathcal E^\vee,\quad a_i\in\mathcal E.\tag*{(1)}\]

% L4–5
Soit $s_0\in\mathcal E$, et posons \deleted{\unsure{Cons|Puis} Posons}

% L4, right under (1)
\[(0\leq i\leq n)\]

% L6, with interline
\[\begin{aligned}s_1&=\tau_0s_0,\ s_2=\tau_1s_1\ (= \tau_1\tau_0s_0),\ \ldots,\ \inserted{s_{i+1}=\tau_i s_i\,\ldots}\\
&s_{n+1}=\tau_n s_n\ (= \tau_n\tau_{n-1}\ldots\tau_0s_0),\end{aligned}\tag*{(2)}\]

% L7, with continuation on right
\[\begin{aligned}\mathcal E_0&=ks_0,\quad\mathcal E_1=k(s_0,\deleted{\unsure{k}}a_0),\ \deleted{\unsure{\mathcal E_2=}}\ldots,\quad\mathcal E_i=k(s_0,\deleted{k}a_0,\ldots\deleted{k}a_{i-1}),\ldots.\\
&\inserted{\mathcal E_{n+1}=k(s_0,\deleted{a_0,}\ldots\deleted{k}a_n).}\end{aligned}\tag*{(3)}\]

% L8, left and below (3)
Soit enfin

% L8
\[\lambda_i=\langle s_0,a'_0\rangle\langle a_0,a'_1\rangle\ldots\langle a_{i-1},a'_i\rangle\quad(0\leq i\leq n)\tag*{(4)}\]

% L8–9, right
\margin{NB $\mathcal E_i$ eng. par $i+1$ \unsure{élts|éls}.}

% L9
\deleted{($\mathcal E_i$ engendré par $i+1$ élt). On a alors}

% L10
1°) On a

% L10–11
\[\begin{aligned}s_i&\in\mathcal E_{\unsure{j|i}}\quad\text{pour }\deleted{\unsure{i,j}}\ \inserted{0\leq i\leq j\leq n+1},\ \text{donc }\deleted{\unsure{\mathcal E}}\ k(s_0,\ldots,s_i)\subset\mathcal E_i\quad(0\leq i\leq n+1)\end{aligned}\tag*{(5)}\]
\[s_i\equiv\deleted{\unsure{s_i}}\lambda_{i-1}a_{i-1}\mod\mathcal E_{i-1}\quad(1\leq i\leq n+1)\quad\text{\uuline{si} on \uuline{suppose}}\tag*{(6)}\]

% L12, right, boxed
\[\boxed{a'_i\mid\mathcal E_{i-1}=0\quad(0\leq i\leq n)}\]

% L12–13
2°) Conditions équivalentes (A) et (B)

% L14–15
A) \inserted{On a} \deleted{Les sections $s_i$ $(0\leq i\leq n+1)$ sont lin. indépendantes,} \deleted{(3) les $(s_i)_{0\leq i\leq n+1}$ \unsure{sont|est} libre}

% L16–18
\[\tau_i\tau_j=\tau_j\tau_i\quad\text{pour }0\leq i,j\leq n,\quad i,j\text{ non consécutifs}\tag*{(\deleted{\unsure{4}}7)}\]
\[\tau_i s_0=s_0\quad\text{pour }1\leq i\leq n\tag*{(\deleted{\unsure{5}}8)}\]
\[(s_i)_{0\leq i\leq n+1}\quad\text{libre dans }\mathcal E\tag*{(\deleted{\unsure{6}}9)}\]

% L16–19, right
% slanted
\margin{il suffit que ce soit vrai \unsure{mod hom|mod hom.}, ou encore dans $\check{\mathbb P}(\mathcal E)$}

% L20
B) On a

% L21–24
\[\langle a_i,a'_j\rangle=0\quad\deleted{\unsure{chaque}}\ \text{pour }0\leq i,j\leq n,\quad i,j\text{ non conséc.}\tag*{(\deleted{\unsure{7}}10)}\]
\[\langle s_0,a'_i\rangle=0\quad\text{pour }1\leq i\leq n\tag*{(\deleted{\unsure{8}}11)}\]
\[(s_0,\deleted{\unsure{a}}a_0,\ldots,a_n)\quad\text{libre dans }\mathcal E.\tag*{(\deleted{\unsure{9}}12)}\]
\[\lambda_n\in k^*\tag*{(13)}\]

% L25
3°) De plus, ces conditions impliquent que

% L26–27
\[\mathcal E_i=k(s_0,\ldots s_{i-1})\qquad0\leq i\leq n+1\tag*{(14)}\]
\[\deleted{\unsure{(14)}\quad s_i\equiv\lambda_{i-1}a_{i-1}\quad\unsure{0\leq i}}\]

% L28–30
\[\begin{aligned}&\deleted{\mathcal E_i\ \text{stable par }\tau_0,\ldots,\tau_{i-1}\quad\unsure{(\quad<i}}\\
&\tau_j(\mathcal E_i)=\mathcal E_i\qquad j\ne i\qquad0\leq i,j\leq n\end{aligned}\tag*{(15)}\]
\[\tau_{\deleted{\unsure{j}}i}(\mathcal E_i)\ne\mathcal E_i\qquad0\leq i\leq n\qquad\text{si }k\ne0\tag*{(16)}\]

% archive sheet number: 96
